\documentclass[11pt]{article}
\usepackage[letterpaper,margin=0.85in]{geometry}
\usepackage[T1]{fontenc}
\usepackage{lmodern,amsmath,amssymb,booktabs,xcolor,fancyhdr}
\usepackage[colorlinks=true,linkcolor=blue,urlcolor=blue]{hyperref}
\definecolor{ink}{HTML}{17344A}
\pagestyle{fancy}\fancyhf{}
\fancyhead[L]{\small\textcolor{ink}{UpstreamDrift / Necromatcher}}
\fancyhead[R]{\small Authored Replay Methods}
\fancyfoot[L]{\small Research Implementation; Historical Qualification Pending}
\fancyfoot[R]{\thepage}
\setlength{\headheight}{14pt}
\setlength{\parindent}{0pt}\setlength{\parskip}{7pt}
\newcommand{\code}[1]{\nolinkurl{#1}}
\title{\textcolor{ink}{Necromatcher Authored Replay\\and Impact Handoff Methods}}
\author{UpstreamDrift Technical Reference\\Epic \#11232; Integration Issue \#11464}
\date{4 October 2026 / Edition 1}
\begin{document}
\maketitle
\begin{center}\rule{0.9\linewidth}{0.6pt}\end{center}
\section*{Purpose and Evidence Boundary}
Necromatcher stores historical image observations, native model hypotheses,
kinematic fits, torque profiles and independently executed authored replays.
This reference describes the implemented public path from a stored replay to
the existing impact and flight pipeline. It supplements the separate camera
and source-scope report. It does not replace that report's retained evidence,
the engineering design manual, or a scientific qualification review.

The source checkpoint for extraction and full-vector consumer behavior in
\href{https://github.com/D-sorganization/UpstreamDrift/pull/11359}{PR \#11359} is:
\begin{center}\small
\code{b059f04e91e5fefd056d8c8514fc2b502815cd89}
\end{center}
At that checkpoint the complete extraction metadata remains in the returned
pipeline result; the trajectory wire alone does not persist the replay receipt.
Portable receipt storage and a normal application action remain open acceptance
criteria. This edition documents implemented mathematics, not proposed behavior.

\section*{Scientific Status}
Both \code{scientific_qualified} and
\code{physical_source_time_qualified} are false. Recorded generalized rates use
authored simulation seconds. Container presentation timestamps identify original
images; they do not establish physical swing duration. No finite differences of
historical video timestamps are used to produce clubhead velocity.

Tiger's reviewed matching domain is the half-open interval $[0,191)$: original
frames 0 through 190. Later released-hand follow-through is excluded. This is a
conservative reviewed interval, not a measured release instant or proof of hand
contact in every retained frame. Hogan's retained domain remains $[0,750)$.
Current historical fits remain research previews; none of the mathematics below
certifies camera calibration, anthropometry, shaft flex, rolling shutter,
clinical range of motion, or historical impact speed.

\newpage
\section*{Inputs and Ownership}
The public workspace facade provides \code{ReplayImpactGeometry},
\code{ReplayImpactSelection} and \code{extract_replay_impact_state}.
Extraction reads through the canonical Library's authenticated-read context.
It compiles no engine merely by importing the module; native evaluation requires
the established SDK-first execution context.

\begin{center}\small
\begin{tabular}{p{0.28\linewidth}p{0.64\linewidth}}\toprule
Declaration & Required Meaning\\\midrule
Replay Identity & Immutable registered authored replay and exact parent hashes.\\
Body and Head Point & Explicit native body and body-local point in metres.\\
Face Axes & Unit, perpendicular body-local normal and up vectors.\\
Effective Mass and MOI & Positive authored impact mass in kg and scalar inertia in kg\,m$^2$.\\
Recorded Sample & Nonnegative integer selecting an existing replay sample.\\
Flight Transform & Proper orthonormal rotation and translation in metres.\\
Assumption Descriptions & Nonempty operator declarations for geometry and selection.\\\bottomrule
\end{tabular}
\end{center}

The selected sample is an authored analysis choice, not automatic contact
detection. The head point is not inferred to be a named frame, the solid centre
of mass, or a measured physical clubhead. Effective impact inertia is not inferred
from silhouette fitting. The typed inputs detach their tuples and reject Boolean,
textual and nonfinite numeric coercions.

\section*{Authenticated Parent Chain}
The replay must have the authored-replay schema and explicitly record independent
execution. Its profile, fit, model and capture must have the expected asset kind,
the same session and the exact stored digest. Compiled native coordinate order,
scalar units, engine identity and fit/model/capture bindings must agree with the
recorded trace. A required derivative capability is checked before coordinate
evaluation. Missing capability fails explicitly.

Extraction rechecks the parent chain before returning; the enclosing authenticated
read also guards source changes. This protects the provenance of an extraction.
It does not independently validate the physical realism of those parents.

\newpage
\section*{Point Velocity From Native Derivatives}
Let $q$ denote the recorded generalized coordinates and $v$ their recorded rates.
For a native body-local point $a_0$, the public marker linearizer provides world
position $p_0(q)$ and coordinate Jacobian $J_0(q)$. The point velocity is
\[
  \dot p_0=J_0(q)v.
\]
Mixed translational and rotational coordinates retain compiled metre/radian
order and corresponding rate units. The implementation does not assume that
all coordinates are angular or replace the native Jacobian with an independent
forward-kinematics implementation.

\section*{Rigid Face Triad and Angular Velocity}
Write the declared body-local normal and up vectors as $n$ and $u$; require
$n^Tn=u^Tu=1$ and $n^Tu=0$. Define $b_1=n$, $b_2=u$ and $b_3=n\times u$.
Three mathematical probes are attached at
\[
  a_i=a_0+L b_i,\qquad L=1\,\mathrm{m},\qquad i=1,2,3.
\]
The probe length establishes units for the derivative construction. It is not
a physical one-metre clubhead. Native marker positions and their Jacobians give
\[
 e_i=\frac{p_i-p_0}{L},\qquad
 \dot e_i=\frac{(J_i-J_0)v}{L}.
\]
For a proper rigid triad, $\dot e_i=\omega\times e_i$. Since the triad is
orthonormal, summing $e_i\times(\omega\times e_i)$ gives $2\omega$, hence
\[
  \boxed{\displaystyle\omega=\frac12\sum_{i=1}^3 e_i\times\dot e_i.}
\]
This obtains all three angular-velocity components from native coordinate
derivatives. Zero angular velocity is not substituted when derivatives are absent.

\section*{Rigid Consistency Checks}
The native public body pose must have a proper rotation. Marker identities,
coordinate order and world positions must agree with that pose and the authored
attachments. Define $E$ with rows $e_i^T$ and $\dot E$ with rows $\dot e_i^T$.
The numerical checks require
\[
 EE^T=I,\qquad \dot E E^T+E\dot E^T=0.
\]
The first condition checks orthonormality; the second checks its instantaneous
preservation. These are numerical rigidity checks, not clinical constraints.

\newpage
\section*{Flight Frame and Consumer Preservation}
The trajectory frame is explicitly \code{flight_xfwd_yleft_zup}: forward,
left and up. With authored world-to-flight rotation $R_f$ and translation $t_f$,
\[
 p_f=R_fp_0+t_f,\quad v_f=R_f\dot p_0,\quad
 \omega_f=R_f\omega,\quad n_f=R_fe_1.
\]
Translation affects position only. The stored loft descriptor is
$\operatorname{atan2}(n_{f,z},\sqrt{n_{f,x}^2+n_{f,y}^2})$ in degrees;
the impact solver receives the full face normal.

The existing trajectory handoff preserves a state that declares the supported
flight frame, including off-axis velocity, angular velocity, face normal and
metadata. An incompatible declared frame is rejected before output. The separate
legacy conversion for undeclared providers retains its existing behavior and
must not be used to collapse an extracted replay into speed and nominal loft.

\section*{Retained Metadata and Interchange}
Extraction records the exact replay, profile, fit, model and capture identities
and hashes; recorded sample index and authored time; initial source-frame
identity; geometry and selection declarations; flight head-point position;
basis probe length; and the explicit clock and qualification flags.
Finite JSON serialization is checked before returning the state.

The flight interchange remains \code{swing_sim.ball_flight_trajectory/1} with
exactly six top-level keys. Its aerodynamic provenance has exactly the model
family, model name and parameter digest. Retained integrator samples must survive
viewer interchange without resampling or re-simulation. Additional historical
provenance requires a separate, versioned companion artifact; adding keys to
the existing flight wire would violate its interchange contract.

\section*{Repeatable Operator Procedure}
\begin{enumerate}
\item Recall the exact authored replay and inspect its parent and clock metadata.
\item Establish the selected engine's SDK-first native execution context.
\item Declare the head geometry, effective inertia, sample and proper flight transform.
\item Extract through the public workspace function. Resolve any admission failure.
\item Pass the extracted state to the existing trajectory coordinator and retain
the returned pipeline result together with the trajectory.
\item Inspect retained samples through the existing viewers. Record the assumptions
and qualification limits whenever sharing the result.
\end{enumerate}

\newpage
\section*{Verification Evidence}
The checkpoint includes an independent synthetic screw-motion fixture that
checks native $Jv$, all angular components, flight rotation and full-vector
consumer preservation. Synthetic kinematics exercise the mathematics; they
are not historical swing measurements. Negative cases cover stale parents,
changed bindings, coordinate order and units, malformed metadata, absent
derivatives, nonrigid probes and mid-read parent changes.

Meaningful failing tests reproduced the old consumer's loss of off-axis
velocity and its acceptance of incompatible declared frames. After repair,
the isolated extraction/consumer/viewer group passed 38 cases. It exercised the
actual impact solver; its flight integration fixture substitutes the optional
Rust integration when the wheel is unavailable. It does not substitute the
impact solver. Cold public imports passed with engine and GUI SDK imports blocked.

The broader checkpoint passed 83 Python cases, 45 React cases and 39 feature
parity checks. These establish covered software behavior only. No new historical
native replay or Library mutation occurred during that checkpoint. Local checks
do not establish remote CI success or scientific release approval.

\section*{Remaining Acceptance Work}
A complete application journey still needs a normal replay-to-impact action,
portable receipt persistence and verified recall through the analysis interfaces.
Historical validation additionally requires justified geometry, time calibration,
camera/anatomy uncertainty, human range-of-motion assessment, contact selection
and improved reprojection quality. Mathematical extraction, polished reporting
and passing synthetic tests do not remove those requirements.

Rolling shutter remains a separate conditional sensitivity hypothesis. A bent
shaft image may also reflect shaft flex, exposure blur, interlacing, film scanning
or broadcast processing. This replay-to-impact path neither estimates a readout
time nor converts shaft curvature into measured angular velocity. Preserve the
original pixels and authored hypothesis whenever comparing alternatives.

\section*{Source References}
\begin{itemize}
\item \href{https://github.com/D-sorganization/UpstreamDrift/issues/11232}{Necromatcher Epic \#11232} and
\href{https://github.com/D-sorganization/UpstreamDrift/issues/11464}{Replay Impact Integration \#11464}.
\item Workspace extraction: \code{necromatcher_impact.py}.
\item Typed declarations: \code{necromatcher_impact_contracts.py}.
\item Existing consumer: \code{trajectory_handoff.py}.
\item Solver pipeline: \code{swing_ball_flight_pipeline.py}.
\item Wire owner: \code{flight_trajectory_export.py}.
\item Extraction evidence: \code{test_necromatcher_impact.py}.
\item Declaration evidence: \code{test_necromatcher_impact_contracts.py}.
\item Consumer evidence: \code{test_shot_trajectory_handoff.py}.
\end{itemize}
\end{document}
